% =====================================================================
% Checking the Blockmodel: A Bayesian Workflow for the Literatures Step
% Caffeine Citation project -- working reference, revision 1
%
% Compile with XeLaTeX or LuaLaTeX for Palatino Linotype.
% =====================================================================
\documentclass[11pt]{article}

\usepackage{amsmath}
\usepackage{amssymb}
\usepackage{bm}
\usepackage{booktabs}
\usepackage{array}
\usepackage{longtable}
\usepackage{enumitem}
\usepackage{ragged2e}
\usepackage{changepage}
\usepackage{titling}
\usepackage{caption}
\captionsetup{labelfont=bf}
\usepackage[margin=1in]{geometry}
\usepackage{underscore}
\usepackage{algorithm}
\usepackage{algpseudocode}
\algrenewcommand\algorithmicrequire{\textbf{Input:}}
\algrenewcommand\algorithmicensure{\textbf{Output:}}

\usepackage{fontspec}
\IfFontExistsTF{Palatino Linotype}{\setmainfont{Palatino Linotype}}{%
  \IfFontExistsTF{TeX Gyre Pagella}{\setmainfont{TeX Gyre Pagella}}{}}

\usepackage[x11names, svgnames, rgb]{xcolor}
\definecolor{gray25}{gray}{0.25}
\definecolor{zred}{rgb}{0.69,0.19,0.19}
\usepackage{hyperref}
\definecolor{navlink}{rgb}{0.10,0.22,0.55}
\hypersetup{
  colorlinks=true, linkcolor=navlink, citecolor=navlink, urlcolor=navlink,
  linktoc=all, bookmarksnumbered=true,
  pdftitle={Checking the Blockmodel: A Bayesian Workflow for the Literatures Step},
  pdfauthor={Jonathan H. Morgan and Sarah Delmar}
}

%% Table conventions (applied selectively, table by table).
%% Double rule under table headers: two 0.5pt rules 1pt apart, pulled 1pt
%% closer to the header above and given 1pt of air below. Use with booktabs:
%%   \toprule ... header ... \doublerule ... body ... \bottomrule
\newcommand{\doublerule}{%
  \specialrule{0.5pt}{-1pt}{1pt}%
  \specialrule{0.5pt}{0pt}{1pt}}
%% Ragged-right paragraph column of a given width: P{2.5cm}
\newcolumntype{P}[1]{>{\raggedright\arraybackslash}p{#1}}
%% \setlength{\tabcolsep}{4.5pt} is set inside the tables that use the
%% convention, not globally.

\newcommand{\code}[1]{\texttt{\small #1}}

\newcommand{\customtitle}[1]{%
  \title{\vspace{-1.5em}#1}
  \pretitle{\begin{center}\LARGE\bfseries}
  \posttitle{\par\end{center}\vskip 12pt}
  \preauthor{\begin{center}\normalsize}
  \postauthor{\par\end{center}\vskip 4pt}
  \predate{\begin{center}\small}
  \postdate{\par\end{center}}
}

\newenvironment{springerabstract}{
  \begin{adjustwidth}{2em}{2em}
  \noindent\textbf{Abstract.}\hspace{0.5em}\justifying
}{
  \end{adjustwidth}
}

\setlength{\droptitle}{0pt}
\customtitle{Checking the Blockmodel\\
\large A Bayesian Workflow for the Literatures Step}
\author{
  \shortstack{
    Jonathan H. Morgan, Ph.D.$^{1}$ \quad Sarah Delmar, Ph.D.$^{2}$
  }
}
\date{
  $^1$\,Duke Network Analysis Center\\
  $^2$\,Department of Psychology, University of Chicago\\
  \today
}

\begin{document}
\maketitle

\begin{springerabstract}
The citation blockmodel returned a consensus partition of 1{,}268 groups for
3{,}557 research communities. Nearly half the communities lacked a stable
assignment across posterior draws, and two runs from the same random seed
returned different hierarchies. When we scored cuts of the fitted hierarchy
with a flat degree-corrected blockmodel, the 405-group cut had the shortest
description among the cuts examined. Before using any cut to estimate how the
effect of topical proximity varies across literatures, we need to determine how
the treatment of citation counts, the hierarchy prior, and the fitting
procedure each contribute to the result. The description-length audit has been run; everything that follows it here
has not. What follows is the plan: three stages of checking -- prior
pushforward, prior predictive, and posterior -- with what each measures, what
domain knowledge it needs, and what each outcome would oblige us to do.
\end{springerabstract}

\tableofcontents
\clearpage

% =====================================================================
\section{The decision this document serves}
% =====================================================================

Model~I -- the community-pair model of \S13.2 of the specification outline,
reproduced as equation (1) below -- cannot be estimated until every research
community carries a label saying which literature it belongs to. A blockmodel was fitted to supply those
labels. It returned a hierarchy thirteen levels deep whose base divides 3{,}557
communities into 1{,}268 groups, left 1{,}771 of those communities without a
stable assignment across posterior draws, and produced different answers on two
runs from the same random seed.

Scoring cuts of that hierarchy under a flat degree-corrected blockmodel gives a
curve with one minimum, at the 405-group cut, and one near-free merge, into 168
groups. Every coarser cut adds between 19 and 33 percent of the shortest description
found. The base of the hierarchy -- the 1{,}268 groups -- scores
80 percent longer than the 405-group cut. (The consensus puts 1{,}268 groups
at the base; the audit, which works from the posterior draws, scores a median
of 1{,}256. \S\ref{sec:audit} reconciles the counts.)

Three candidate explanations remain open, and the workflow in
\S\ref{sec:pushforward} to \S\ref{sec:posterior} exists to separate them: the
treatment of citation counts, which fixes how much the model believes it has
seen; the hierarchy prior, which fixes how cheaply fine partitions can be
described; and the fitting procedure itself, whose non-reproducibility is
unexplained. The decision that waits on all three is which cut of the hierarchy
becomes the literature label $\ell$ in Model~I.

% =====================================================================
\section{What a literature is, and why the model needs one}
% =====================================================================

Model~I asks a question about pairs of research communities: when two
communities work on similar questions, do they draw on the same earlier
literature more often than two communities that do not? Shared citation is the
observable, topical proximity is the predictor, and the coefficient linking them
is the quantity of interest.

That coefficient is not allowed to be a single number. It varies two ways. It
drifts across the six era pairs, which asks whether the field's tendency to
connect similar work has changed over fifty years. And it varies across
\emph{literatures}, which asks whether the tendency works the same way
everywhere in the field or differently in different parts of it. Writing
$S_{ijp}$ for the topical proximity of communities $i$ and $j$ in era pair $p$,

\begin{equation}
\log \lambda_{ijp} = \alpha^{K}_{p} + \log E^{K}_{ijp}
  + \underbrace{\alpha_{S,p}\, S_{ijp}}_{\text{drifts over time}}
  + \underbrace{\alpha_{S,\ell}\, S_{ijp}}_{\text{varies by literature}}
  + \bm{z}_{ijp}^{\top}\bm{\gamma} + \theta^{K}_{i} + \psi^{K}_{j}.
\end{equation}

The superscript $K$ marks the coupling model throughout, distinguishing its
era baselines, exposures, and community effects from those of the citation
model that shares the specification. $\theta^{K}_{i}$ and $\psi^{K}_{j}$ are
random intercepts for the citing and cited communities.

The second of those terms needs an $\ell$: a label on every community saying
which literature it belongs to. Nothing else in the project supplies one.
Journals do not, because a cognitive paper can appear in a pharmacology journal.
Keywords do not, because they describe subject matter rather than position.
Communities themselves do not, because they are found separately within each
seven-year era and carry no identity across eras.

A literature, in the sense the model needs, is a group of communities that
connect to the rest of the field in the same way. Two pharmacokinetics
communities in adjacent eras may barely cite one another and still belong
together, because the rest of the field draws on them interchangeably. That
criterion -- sameness of position rather than density of contact -- is what a
blockmodel estimates, and it is why a blockmodel is in the pipeline at all.

The specific instrument is a \emph{nested degree-corrected stochastic
blockmodel}, fitted with the \code{graph-tool} library. Three words in that name
do real work. \emph{Stochastic blockmodel}: communities are sorted into groups
so that the citation traffic between any two groups is described by a single
rate, and the sorting that describes the traffic most economically wins.
\emph{Degree-corrected}: large communities have more connections by virtue of
being large, and the model accounts for that so the groups reflect affinity
rather than volume. \emph{Nested}: the groups are themselves grouped, and those
groups grouped again, so the result is a hierarchy rather than a single
partition.

% =====================================================================
\section{What the fit returned}
% =====================================================================

The network handed to the blockmodel has one node per community -- 3{,}557 of
them, across eras 11 to 17 -- and a directed edge wherever one community's works
cite another's, weighted by the number of citations. It holds 36{,}105 distinct
community pairs and 583{,}955 citations, with pair weights running from 1 to
4{,}646.

Table~\ref{tab:symptoms} records what came back. Each row is a result that would
be unremarkable on its own and is not unremarkable in company.

\begin{table}[!t]
\centering
\setlength{\tabcolsep}{4.5pt}
\caption{What the fit returned, and why each result is a problem}
\label{tab:symptoms}
\begin{tabular}{@{}P{3.6cm}P{3.2cm}P{7.4cm}@{}}
\toprule
\textbf{Result} & \textbf{Value} & \textbf{Why it is a problem}\\
\doublerule
Groups at the base of the hierarchy
  & 1{,}268
  & 3{,}557 communities in 1{,}268 groups is 2.8 communities apiece. Most
    groups would then contribute a handful of pairs over a narrow range of
    topical proximity, leaving $\sigma_{\ell}$ estimated largely from its
    prior. See \S\ref{sec:downstream}.\\
Groups holding one community
  & 690 (54\%)
  & Half the groups are not groups. The communities in them are not marginal
    either: their median degree is 34.5 against 5 for the rest, and their median
    size is 89 works against 4.\\
Communities with no settled assignment
  & 1{,}771 (49.8\%)
  & Across posterior draws, half the communities fail to place in any one group
    on 90 percent of draws. The consensus partition is an average of
    disagreeing answers.\\
Consensus against point estimate
  & 1{,}268 vs 846
  & The consensus is \emph{finer} than the fit it summarizes. Averaging
    partitions that disagree produces their intersection, and intersections
    splinter.\\
Reproducibility
  & two answers
  & Two runs on identical input with the same seed returned different
    hierarchies, 2{,}451 nats -- natural-log units of information -- apart in
    description length.\\
Compute
  & 15--20 hours
  & Of which 7.6 hours was reconciling draws, a step whose cost scales with
    the number of groups being matched. 499 were collected, every second one
    aligned, so 250 entered the consensus.\\
\bottomrule
\end{tabular}
\end{table}

% =====================================================================
\section{The audit that turned a suspicion into a measurement}
\label{sec:audit}
% =====================================================================

Three group counts appear in this document and they are three different
objects. The \emph{point estimate} from the description-length heuristic has
846 groups at its base. The \emph{consensus}, formed by reconciling 250 of the
499 posterior draws, has 1{,}268. The \emph{median across the 50 draws} scored in
the audit below has 1{,}256. Nothing is inconsistent here; a consensus of
disagreeing partitions tends toward their intersection, which is why it is
finer than the estimate it summarizes.

A hierarchy offers a choice of granularity, and the obvious question is which
level to treat as the literature. The hierarchy itself cannot answer: it runs
1{,}268, 405, 168, 82, 44, 25, 14, 8, 4, 2, 1, with successive ratios of 3.1,
2.4, 2.0, 1.9, and then 1.8 to 2.0 the rest of the way down. The base thins
fastest and the lower levels are steady; nowhere is there a plateau or a
gap. A dendrogram
drawn from it is no help either, because its horizontal axis is level number
rather than any measure of dissimilarity -- the branches have no lengths, so
nothing in the picture marks a natural cut.

Description length supplies the lengths. Take each level's grouping, project it
back onto the 3{,}557 communities, hold it fixed, and ask what a flat
degree-corrected blockmodel with that grouping costs to describe. Coarsening
always loses fit and always saves complexity; the level where the loss begins to
outrun the saving is the granularity the network supports.

\begin{table}[!t]
\centering
\setlength{\tabcolsep}{4.5pt}
\caption{Flat-model description length of each cut of the fitted hierarchy,
  median over 50 posterior draws. The three degenerate top levels, all of one
  group, are omitted.}
\label{tab:entropy}
\begin{tabular}{@{}rrrrP{5.4cm}@{}}
\toprule
\textbf{Groups} & \textbf{DL (nats)} & \textbf{Merge cost} & \textbf{vs best}
  & \textbf{Reading}\\
\doublerule
1{,}256 & 1{,}327{,}232 & --- & $+80.1\%$
  & base of the hierarchy\\
405 & 736{,}908 & $-590{,}324$ & --- & lowest-scoring cut examined\\
168 & 750{,}206 & $+13{,}298$ & $+1.8\%$
  & the only near-free merge\\
82 & 971{,}404 & $+221{,}198$ & $+31.8\%$ & \\
44 & 1{,}120{,}231 & $+148{,}827$ & $+52.0\%$ & \\
25 & 1{,}256{,}995 & $+136{,}764$ & $+70.6\%$
  & assignment confidence stops improving; no singletons\\
14 & 1{,}463{,}906 & $+206{,}911$ & $+98.7\%$ & \\
8 & 1{,}646{,}272 & $+182{,}366$ & $+123.4\%$
  & two groups hold 2 percent between them\\
4 & 1{,}868{,}082 & $+221{,}810$ & $+153.5\%$ & \\
2 & 2{,}110{,}950 & $+242{,}868$ & $+186.5\%$ & \\
\bottomrule
\end{tabular}
\end{table}

Two readings follow, and they point in opposite directions. The base of the
hierarchy scores \emph{80 percent longer} under the flat model than the
405-group cut does. Whatever produced 1{,}268 groups, it was not a flat
accounting of the citation record. But coarsening scores badly too: every merge below 168 groups adds between 19
and 33 percent of the shortest description found, so a reading of eight
literatures gives up more than the whole of that description.

Two cautions on reading Table~\ref{tab:entropy}. The comparison is among
\emph{cuts of the fitted hierarchy}, scored under a flat model; no flat
partition was optimized independently, so 405 is the lowest-scoring cut examined
rather than the best flat partition available. And the scoring model is not the
model that produced the hierarchy -- the nested fit's own description length is
about 225{,}000 nats, far below anything in the table, because the hierarchical
prior of \S\ref{sec:model} charges much less for the same granularity. The
table prices cuts on a common scale; it does not adjudicate between the two
models.

The fifty posterior draws lie almost on top of one another in this curve. The
shape is not one fit's idiosyncrasy.

The audit also recorded three facts about particular cuts that description
length does not capture, and \S\ref{sec:downstream} leans on them. At the
25-group cut, assignment confidence -- the share of draws placing a community
in its modal group -- stops improving as the cut coarsens; 24 of the 25 groups
span more than one era; and no group is a singleton. At the 8-group cut, two
of the eight groups hold 2 percent of the communities between them, so a
reading of eight literatures is in practice a reading of six.

% =====================================================================
\section{The model being fitted}
\label{sec:model}
% =====================================================================

The assumptions at issue live in particular terms of this model, so it is set
out here rather than named. The account is Peixoto's, as implemented in
\code{graph-tool}; the notation is his.

\subsection{One partition}

Write $N = 3{,}557$ for the number of communities and $\bm{A}$ for the matrix of
citation counts, so that $A_{ij}$ is the number of citations from community
$i$'s works to community $j$'s. A partition $\bm{b}$ assigns every community a
group, $b_i \in \{1, \dots, B\}$. Two quantities summarize what the partition
sees: the citations running between groups,
\[
e_{rs} \;=\; \sum_{ij} A_{ij}\,\delta_{b_i r}\,\delta_{b_j s},
\]
with margins $e^{\text{out}}_{r} = \sum_s e_{rs}$ and $e^{\text{in}}_{s} =
\sum_r e_{rs}$; and the degrees $k^{\text{out}}_{i}$ and $k^{\text{in}}_{i}$ of
each community.

The generative story is that citations arrive as independent Poisson draws
whose rate depends on the groups at either end and on how active each community
is:
\begin{equation}
A_{ij} \sim \mathrm{Poisson}\!\left(
  \theta^{\text{out}}_{i}\,\theta^{\text{in}}_{j}\,\lambda_{b_i b_j}\right).
\end{equation}
Here $\lambda_{rs}$ is the affinity between groups $r$ and $s$, and
$\theta^{\text{out}}_{i}$, $\theta^{\text{in}}_{i}$ are the propensities of the
individual community. Those $\theta$ terms are the degree correction. Without
them a community that cites heavily would be grouped with other heavy citers
whatever it cites, and with community sizes in this corpus running from 2 works
to 2{,}769, that failure mode is not hypothetical.

Maximizing over $\bm{\theta}$ and $\bm{\lambda}$ leaves a quantity that depends
only on the partition,
\begin{equation}
\mathcal{S}(\bm{b}) \;=\; -\sum_{rs} e_{rs}\,
  \ln\!\frac{e_{rs}}{e^{\text{out}}_{r} e^{\text{in}}_{s}},
\label{eq:entropy}
\end{equation}
which is how many nats -- units of information, like bits but measured in
natural logarithms -- it takes to transmit the citation matrix once the
receiver knows the partition. Smaller is better, and $\mathcal{S}$ falls
monotonically as groups are added -- a partition with one community per group
describes the data perfectly and explains nothing.

\subsection{What stops it from using every group it can}

The brake is the cost of transmitting the model itself. The full description
length adds to \eqref{eq:entropy} the cost of the partition, the between-group
counts, and the degrees:
\begin{equation}
\Sigma \;=\; \underbrace{\mathcal{S}(\bm{b})}_{\text{the data}}
  \;+\; \underbrace{\mathcal{L}}_{\text{the model}},
\qquad
\mathcal{L} = -\ln P(\bm{e}) - \ln P(\bm{k} \mid \bm{e}, \bm{b})
  - \ln P(\bm{b}).
\label{eq:dl}
\end{equation}
Minimizing $\Sigma$ is identical to maximizing the posterior
$P(\bm{b} \mid \bm{A})$ under those priors. The three terms of $\mathcal{L}$
\emph{are} the priors, and they are these:
\begin{align}
P(\bm{b}) &= \frac{\prod_r n_r!}{N!}
  \,\binom{N-1}{B-1}^{\!-1}\,\frac{1}{N},
  &&\text{group sizes $n_r$, uniform given $B$,}\label{eq:pb}\\[2pt]
P(\bm{e}) &= \left(\!\!\binom{B^2}{E}\!\!\right)^{\!-1},
  &&\text{uniform over count matrices,}\label{eq:pe}\\[2pt]
P(\bm{k} \mid \bm{e}, \bm{b}) &= \prod_r
  \left(\!\!\binom{n_r}{e^{\text{out}}_{r}}\!\!\right)^{\!-1}
  \left(\!\!\binom{n_r}{e^{\text{in}}_{r}}\!\!\right)^{\!-1},
  &&\text{uniform over degree sequences,}\label{eq:pk}
\end{align}
where $\left(\!\binom{n}{k}\!\right) = \binom{n+k-1}{k}$ counts the ways to
choose $k$ items from $n$ types with repetition, and $E$ is the total number of
edges. The directed model carries separate in- and out-degree terms, as shown.

One caution on equation~\eqref{eq:pk}. It is the \emph{uniform} degree prior, which
\code{graph-tool} offers as an option; the library's default is
\code{degree_dl_kind="distributed"}, a non-uniform hyperprior that models the
degree distribution itself. Which was in force during our fit has not been
checked, and it should be, because a hyperprior on degrees is one of the few
places in this model where something resembling domain knowledge could enter.

What the three terms have in common is that each is uniform over some
combinatorial set, and none of them mentions citation, or science, or anything
else about where the data came from.

\subsection{From one partition to a hierarchy}

The flat prior $P(\bm{e})$ above has a known and unwanted consequence: it
cannot resolve more than about $\sqrt{E}$ groups, because the cost of
describing a $B \times B$ count matrix grows faster than the fit improves.
The ceiling is a function of how many edges the model believes it has, which
is why \S\ref{sec:weights} matters here too: at the 36{,}105 distinct pairs
it is on the order of 190, and at the $E = 583{,}955$ the multiplicity
treatment declares it is on the order of 760. The constant in front of
$\sqrt{E}$ depends on the average degree, so these are magnitudes rather than
values. Peixoto states the limit in the familiar form $O(\sqrt{N})$, which
coincides with $\sqrt{E}$ only for sparse graphs with $E = O(N)$; under the
multiplicity treatment ours is not sparse in that sense, and $\sqrt{N} = 60$
would be off by an order of magnitude. Either way it is a detection limit
-- a statement about what
the flat model can resolve, not about how many literatures the field
contains.

The nested model removes the ceiling by refusing to describe the count matrix
flatly. Instead it treats $\bm{e}$ as a network in its own right -- one with
$B_1$ nodes -- and partitions \emph{that}, recursively, until one group remains:
\begin{equation}
P(\{\bm{e}_l\}, \{\bm{b}_l\}) \;=\;
  \prod_{l=1}^{L} P(\bm{e}_l \mid \bm{e}_{l+1}, \bm{b}_{l})\,P(\bm{b}_{l}),
  \qquad B_L = 1 .
\label{eq:nested}
\end{equation}
Our fit has $L = 13$ levels, with $B_1 = 1{,}268$ at the base. The top three
collapse to a single group and carry no information, which is why
Table~\ref{tab:entropy} lists ten cuts rather than thirteen.

Equation~\eqref{eq:nested} is the single most consequential line in this
document. It is the reason a grouping into 1{,}268 blocks is cheap to transmit:
the blocks need not be described one by one against a flat background, only
against the coarser grouping above them. Note where this assumption sits. It
is a prior over the count matrices, not over the partition -- the marginal
prior on the number of groups is uniform, as \S\ref{sec:pushforward} shows --
so it reaches $\bm{b}$ through the marginal likelihood. That is why the
instrument for probing it is the prior predictive and not a pushforward.
Equation~\eqref{eq:nested} also fixes the prior over depth: each $B_l$ is
drawn uniformly from $\{1, \dots, B_{l-1}\}$, so the hierarchy is a chain
that descends from $N$ to one group in a random number of steps.

\subsection{Weights, and the quantity $E$}
\label{sec:weights}

Citation counts enter through $E$, and $E$ appears in $P(\bm{e})$. Passing the
counts as \code{eweight} declares them to be edge multiplicities: a pair joined
by 4{,}646 citations contributes 4{,}646 to $E$, so $E = 583{,}955$ rather
than the 36{,}105 distinct community pairs. Since $P(\bm{e})$ is a multiset
coefficient over $B^2$ cells and $E$ draws, a larger $E$ supports a larger $B$
at the same cost. The likelihood choice and the prior are therefore coupled
through a single integer, which is why \S\ref{sec:order} settles the weights
before checking anything else.

\subsection{How it was fitted, and with what}

The implementation is \code{graph-tool} 3.9, Tiago Peixoto's C\texttt{++}
library with Python bindings, installed in its own \code{conda} environment at
\code{/opt/gt-env} and invoked by absolute path. The version string and the
argument names below were read from the installed library rather than from
documentation: an earlier script passed \code{state_args}, the name used in
the 2.x series, and 3.9 warned of \code{unrecognized keyword arguments} and
discarded it. Every call here is written against signatures confirmed
by inspection. Four calls do the work:

\begin{itemize}[leftmargin=1.5em, itemsep=0.3em]
\item \code{minimize_nested_blockmodel_dl} -- an agglomerative heuristic that
  searches for the partition minimizing \eqref{eq:dl}, returning a point
  estimate in about 45 seconds. Degree correction and the weights enter here,
  through \code{base_state_args}.
\item \code{mcmc_equilibrate}, with merge-and-split proposals -- Markov chain
  Monte Carlo over partitions. Starting from the heuristic's answer it found a
  description length 18{,}329 nats lower, so the heuristic alone was well short
  of the mode.
\item \code{PartitionModeState} -- group labels are arbitrary and change
  between draws, so nothing can be averaged until the draws are put on a common
  labeling. Of the 499 draws collected, every second one was aligned; those
  250 took 7.6 hours.
\item \code{get_max_nested} and \code{get_marginal} -- the consensus partition
  and, for each community, the share of draws placing it in each group.
\end{itemize}

\code{graph-tool} exposes arguments for the likelihood -- degree correction on
or off, weights as multiplicities or as covariates, nested or flat. It exposes
no argument for $P(\bm{b})$. The priors in equations~\eqref{eq:pb} to
\eqref{eq:pk} are the ones we get.

% =====================================================================
\section{Where the assumptions are}
% =====================================================================

Two choices shaped everything above, and neither was examined before the fit.

\subsection{The likelihood: citation counts as multiplicities}

Section~\ref{sec:weights} gave the mechanism; what matters here is that it was
never examined. \code{eweight} declares the counts to be edge multiplicities,
so the pair joined by 4{,}646 citations enters the likelihood as 4{,}646
parallel edges rather than as one relationship of high intensity. That single
argument fixes $E$, $E$ fixes $P(\bm{e})$, and $P(\bm{e})$ is one of the three
terms that decides how many groups the model will tolerate.

The alternative treats the count as a covariate attached to one edge, through
\code{recs} with a discrete Poisson record type. That is a different generative
story about the same observations, and the one usually preferred when weights
are heavily skewed. Neither is obviously right. What is wrong is having
chosen between them in a single unremarked line.

\subsection{The prior: reference priors with no domain content}

\code{graph-tool}'s priors over partitions are uniform within microcanonical
ensembles -- every arrangement consistent with a fixed set of totals gets equal
weight -- chosen so that the description length is a clean coding argument with
nothing tuned. They encode nothing whatever about citation networks. Nothing in
them says that in-degree is heavy-tailed, that reference lists are bounded, that
citation runs backwards in time and almost never forwards, or that a field of
this size holds tens of specialties rather than thousands.

Being empty of domain content does not make them weak in effect. The nested
prior exists precisely to make deep hierarchies with many small blocks cheap to
describe, and Table~\ref{tab:entropy} measures how much that costs once the
subsidy is withdrawn. We therefore have a prior with no knowledge in it and
large consequences out of it: not uninformative where it matters, and not
informed where it should be.

There is a domain-informed prior available, and it has been stated in
conversation without being written down: a field of this size, over fifty years,
holds something like twenty to thirty literatures. That is a weakly informative
claim, defensible from the co-citation literature on specialty structure, and
the model had no access to it. Formalizing it is the business of
\S\ref{sec:pushforward}.

\subsection{Why the order matters}

The two cannot be checked independently. Informativeness is a property of the
prior and the likelihood together, not of the prior alone -- a flat prior on a
logistic regression coefficient is wildly informative about outcomes. Here, the
same prior over partitions implies one predictive distribution when the model
believes it has 583{,}955 observations and quite another when it believes it has
36{,}105. Settling the likelihood therefore comes first, and every check below
assumes it has been settled.

% =====================================================================
\section{Stage one: prior pushforward checks}
\label{sec:pushforward}
% =====================================================================

A pushforward check pushes the prior through a \emph{function of the
parameters} and inspects the distribution that comes out. It stops short of
generating data, which makes it cheap, and it is the natural first instrument
whenever the quantity in dispute is itself a parameter.

The quantity in dispute is the number of groups. Nobody has intuitions about a
prior over partitions of 3{,}557 objects; everybody has intuitions about how
many literatures a field contains. The pushforward is supposed to convert the
first into the second.

\subsection{What the prior over $B$ actually is}

It can be read off equation~\eqref{eq:pb} without sampling anything. The
partition prior factors as $P(\bm{b}) = P(\bm{b} \mid \bm{n})\,P(\bm{n} \mid
B)\,P(B)$ with $P(B) = 1/N$, and the remaining two terms of $\mathcal{L}$ are
normalized conditionals -- $P(\bm{e})$ given $B$ and $E$, and $P(\bm{k})$
given $\bm{e}$ and $\bm{b}$ -- so summing them out leaves $P(\bm{b})$
unchanged. The marginal prior over the number of groups is therefore
\textbf{uniform on $\{1, \dots, 3{,}557\}$}, with mean $1{,}779$. The same
holds at every level: $B_l$ is uniform on $\{1, \dots, B_{l-1}\}$.

Three things follow, and together they redirect this stage.

\emph{The pushforward on $B$ cannot discriminate.} A flat distribution has no
mass to concentrate in the tens or near a thousand, so a question posed as
``where does the prior put $B$?'' is answered analytically and tells us
nothing about the fit.

\emph{The count is the likelihood's choice.} The prior has no preference; the
heuristic returned 846 and the consensus 1{,}268 because the marginal
likelihood selected them. A fine partition therefore cannot be read as the
prior's doing, and nothing about the base count was settled before the
citations were read.

\emph{The expensive assumption is elsewhere.} What makes a fine partition cheap
is the nested prior over the count matrices, equation~\eqref{eq:nested}: a
$B \times B$ matrix need not be described against a flat background, only
against the coarser grouping above it. That is a prior over
\emph{parameters}, and it acts on $\bm{b}$ only through the marginal
likelihood. A pushforward on $B$ is the wrong instrument for it. The right one
is the prior predictive of \S\ref{sec:predictive}, and specifically its
flat-against-nested arm.

\subsection{What the pushforward is still good for, and why it needs no
  simulation}

Two quantities are not uniform and are worth having. Both come from
$P(\bm{b})$ alone -- neither touches $\bm{e}$ or $E$ -- and both have
closed forms, so Stage one is arithmetic rather than sampling. The figures
below were computed by the recursion in Algorithm~\ref{alg:pushforward} and
should be reproduced before they are relied on.

\emph{The share of singleton groups.} $P(\bm{n} \mid B)$ is uniform over
compositions of $N$ into $B$ positive parts, and under that law the
probability that a given group has exactly one member is
\begin{equation}
P(n_r = 1 \mid B) \;=\; \frac{\binom{N-2}{B-2}}{\binom{N-1}{B-1}}
  \;=\; \frac{B-1}{N-1},
\label{eq:singleton}
\end{equation}
which by exchangeability is also the expected share of singleton groups.
Conditional on $B = 1{,}268$ the prior expects $35.6$ percent singletons;
conditional on $B = 846$, $23.8$ percent. The fit's 54 percent exceeds both.
The marginal figure, averaged over the uniform $B$, is exactly 50 percent,
which is close to the fit -- but that closeness is an artifact of the uniform
$B$, half of whose mass sits above 1{,}779 groups where nearly every group is
a singleton. The conditional figures are the ones that bear on the fit, and
they say the prior did not anticipate this many singletons.

\emph{Hierarchy depth.} Write $d$ for the number of levels until the chain
$B_0 = N,\; B_l \sim \mathrm{Uniform}\{1, \dots, B_{l-1}\}$ first reaches
one group. A level may repeat the count above it, since the uniform draw
includes $B_{l-1}$ itself, and such a level still counts toward the depth.
The distribution of $d$ is a one-dimensional recursion over $B$, exact and
instantaneous. At $N = 3{,}557$ the mean depth is $9.75$ levels, and the
probability of eleven or more is $0.38$. Our fit descends through ten
non-degenerate levels to a single group at the eleventh, so a hierarchy of
this depth is ordinary under the prior rather than extreme. Conditioning on
the fitted base count moves the mean to $9.72$ given $B_1 = 1{,}268$ and to
$9.32$ given $B_1 = 846$; neither makes eleven levels remarkable.

\emph{Some level in the elicited range.} This one needs no recursion. For a
band $[\mathrm{lo}, \mathrm{hi}]$ and any starting count above it, the
probability that the chain passes through a level inside the band is
\begin{equation}
P\big(\text{some } B_l \in [\mathrm{lo}, \mathrm{hi}]\big)
  \;=\; \frac{\mathrm{hi} - \mathrm{lo} + 1}{\mathrm{hi}},
\label{eq:band}
\end{equation}
exactly, independent of $N$ and of $B_1$. The identity follows from
self-consistency: write $v(b)$ for the probability of reaching the band from
$b > \mathrm{hi}$; the next count is uniform on $\{1, \dots, b\}$, so
$v(b) = \frac{1}{b}\big[v(b) + \sum_{b' < b} v(b')\big]$, and the constant
$v = (\mathrm{hi} - \mathrm{lo} + 1)/\mathrm{hi}$ solves this at every $b$
above the band once it holds at $b = \mathrm{hi} + 1$, where it is immediate.
For twenty to thirty groups the figure is $11/30 = 0.367$; for the wider band
of ten to forty it is $31/40 = 0.775$. The prior gives a readable level with
probability a little over one in three, whatever the fit did, which is neither
the near-certainty the domain expectation would want nor a rejection of it.

\subsection{The domain threshold}

The threshold cannot be stated as ``$P(B_{\ell} > 100) \approx 0.01$ at the
cut intended for $\ell$,'' because which cut becomes $\ell$ is the output of
this whole process; a threshold defined by reference to its own result is
circular. Two non-circular forms are available. Fix the level in advance --
declare that $\ell$ will be read from level five, say -- and the threshold
applies to a well-defined quantity. Or state it as a property of the
hierarchy: the prior should place substantial mass on hierarchies possessing
\emph{some} level with between twenty and thirty groups. The second is closer
to what the domain knowledge actually supports, since the co-citation
literature speaks to how many specialties a field has and not to the depth at
which a blockmodel should find them; and it is the form that
equation~\eqref{eq:band} answers in closed form.

\subsection{What \code{graph-tool} does and does not expose}

The priors in equations~\eqref{eq:pb} to \eqref{eq:pk} are written down, so
they can be sampled directly in a dozen lines of either language. What cannot
be replaced from outside the library is the \emph{fitting}: the samplers and
the heuristic are built against those priors internally, and there is no
argument for supplying a different $P(\bm{b})$.

Two partial exceptions should be tried before concluding that nothing is
available. \code{B_min} and \code{B_max} constrain the number of base groups
directly. \code{clabel} and \code{pclabel} constrain which vertices may share a
group. Both are crude next to a stated prior -- a hard interval is not a
distribution -- and \code{B_max} has a further defect: it constrains the
\emph{base}, while the domain expectation is about some upper level. A model
can hold a thousand base blocks and still deliver twenty-five literatures five
levels up, so capping the base at thirty asks for something the expectation
never demanded. There is no argument for capping an upper level. \code{B_max}
is therefore used here as a proxy that is known to be too strong, and the
reading of its result is set accordingly: the quantity recorded is the
nested model's own description length under the cap against its unconstrained
value, and a penalty on the scale of the near-free merge in
Table~\ref{tab:entropy} -- a few percent -- would count as acceptable. A
penalty on the scale of that table's coarser cuts, tens of percent, would
not, but because the cap is too strong a large penalty does not by itself
show that the prior is unrevisable; it shows that this argument cannot do the
revising.

\begin{algorithm}[!t]
\caption{Prior pushforward. Every quantity here follows from
  equation~\eqref{eq:pb} and the level recursion of
  equation~\eqref{eq:nested}; nothing is sampled and $E$ does not enter.}
\label{alg:pushforward}
\begin{algorithmic}[1]
\Require $N = 3{,}557$; the fitted $B_1$ (846 or 1{,}268) and depth (eleven
  levels to one group); the elicited band $[20, 30]$
\Statex \textbf{Number of groups}
\State $B_1 \sim \mathrm{Uniform}\{1, \dots, N\}$, mean $1{,}779$
  \Comment{report; it carries no information about the fit}
\Statex \textbf{Singleton share, by equation~\eqref{eq:singleton}}
\State $z(B) \gets (B-1)/(N-1)$ for $B \in \{846,\ 1{,}268\}$
\Comment{35.6 and 23.8 percent; marginal over $B$ is 50 percent}
\Statex \textbf{Depth, by recursion over $B$}
\State $p_0 \gets \delta_N$ \Comment{or $\delta_{B_1}$ to condition on the fit}
\For{$l = 1, 2, \dots$ until the mass at $B = 1$ exceeds $1 - 10^{-6}$}
  \State $p_l(b) \gets \sum_{b' \ge b} p_{l-1}(b')/b'$ for $b \ge 2$;
    mass newly arriving at $b = 1$ is $P(d = l)$
  \Comment{the sum includes $b' = b$: a level may repeat the count above it}
\EndFor
\Statex \textbf{A level in the elicited band, by equation~\eqref{eq:band}}
\State $P(\text{some } B_l \in [20, 30]) \gets 11/30$
  \Comment{exact; the same for every $N$ and $B_1$ above the band}
\Ensure $P(d)$ and its mean; $P(d \ge 11)$; $z(B_1)$; $11/30$
\State \textbf{Decisions:} is 54 percent singletons above $z(B_1)$? is
  $P(d \ge 11)$ small? is $11/30$ an acceptable chance of a readable level?
\State \textbf{Then:} rerun the heuristic under \code{B_max} $= 30$ and record
  the nested description length against the unconstrained fit, reading the
  result as \S\ref{sec:pushforward} directs
\end{algorithmic}
\end{algorithm}

\begin{table}[!t]
\centering
\setlength{\tabcolsep}{4.5pt}
\caption{Outcomes of Stage one and what each requires. The rows concern depth
  and group sizes; the number of groups is uniform under the prior and is not
  at issue here.}
\label{tab:pushforward}
\begin{tabular}{@{}P{4.4cm}P{4.2cm}P{5.6cm}@{}}
\toprule
\textbf{If the prior} & \textbf{Then} & \textbf{And we}\\
\doublerule
expects fewer singletons than the fit returned, and depths like the fit's
  & the singletons are the data or the likelihood talking; the depth is not
    evidence of anything
  & carry the singleton question to Stage two, where the nested prior over
    $\bm{e}$ can be isolated\\
expects at least as many singletons as the fit returned, conditional on
  $B_1$
  & that part of the fit's shape was anticipated by the prior, independent of
    any citation
  & report it, and stop treating the base-level group sizes as a finding
    about the field\\
gives a level in the elicited band with probability $11/30$, which the
  domain expectation finds too low, and \code{B_max} cannot reach the band
  without a penalty of tens of percent
  & the available arguments cannot bring the model into line with domain
    knowledge
  & reimplement the degree-corrected blockmodel with a hand-written sampler and
    a stated prior\\
\bottomrule
\end{tabular}
\end{table}

On the figures above, the first row is the one the arithmetic already
selects: 54 percent singletons against an expected 36, and a depth that the
prior finds unremarkable. The third row remains open until the \code{B_max}
rerun has been priced. A workflow exists to make assumptions visible and
revisable; if the only revisions available are a hard interval on $B$ and a
constraint on which vertices may share a group, then the prior is not
revisable in the sense the workflow requires, and demonstrating that is more
useful than suspecting it.

% =====================================================================
\section{Stage two: prior predictive checks}
\label{sec:predictive}
% =====================================================================

A predictive check pushes the prior all the way through the data-generating
process and asks whether the networks it produces could plausibly be citation
networks. Where the pushforward asks whether the implied shape of the
hierarchy is sane, the predictive check asks whether the implied
\emph{universe of networks} contains anything resembling ours. This is also
the first stage at which $E$ matters, since the count matrices are drawn
here and not before.

\paragraph{Procedure.} Draw a partition from the prior, draw the between-group
counts and the degrees from theirs, generate a network, and compute its
summary statistics. Repeat a few hundred times. Then overlay the observed
values. The draws are \emph{ancestral}: each quantity is drawn in the order the
prior generates it, conditional on the one before.

\paragraph{What to compare, and against what.} The statistics worth checking are
those the model does not fit by construction. Degree distribution is a poor
check, because degree correction reproduces it by design. Table~\ref{tab:stats}
lists the useful ones with their sources.

\begin{table}[!t]
\centering
\setlength{\tabcolsep}{4.5pt}
\caption{Statistics for the prior predictive comparison. Forward-in-time
  citation is not among them: 37 of the 1{,}098{,}984 arcs between
  analytic-era works run forwards, but that is a property of the article
  network, and the community model has no time in it to reproduce.}
\label{tab:stats}
\begin{tabular}{@{}P{3.1cm}P{4.6cm}P{6.4cm}@{}}
\toprule
\textbf{Statistic} & \textbf{Observed in our network}
  & \textbf{Domain expectation and source}\\
\doublerule
Reciprocal community pairs
  & to be counted: pairs $(r,s)$ with $e_{rs} > 0$ and $e_{sr} > 0$, out of the
    36{,}105 pairs in the fitted graph
  & Low but not zero. Communities in adjacent eras can cite each other in both
    directions even though individual works cannot. Distinct from the
    article-level forward-citation count noted in the caption.\\
In-degree tail
  & degree 2 to 522, median 5
  & Heavy-tailed, from cumulative advantage -- Price's account of why
    well-cited papers attract further citation.\\
Weight skew
  & pair weights 1 to 4{,}646
  & Community-pair intensity inherits the same skew. A prior implying
    homogeneous weights is at odds with the data by three orders of
    magnitude.\\
$k$-core profile
  & coreness 2 to 74, with 125 communities at 74 and no vertex below 2
  & A dense core with a connected periphery. This is the one statistic that
    nothing in the model fits, which makes it the sharpest check available.\\
Group-size distribution
  & 54 percent of fitted groups hold one community
  & Specialty sizes are skewed but not dominated by singletons. A prior that
    expects half its groups to be singletons expects something unlike a
    field.\\
\bottomrule
\end{tabular}
\end{table}

\begin{algorithm}[!t]
\caption{Prior predictive check, simulating the hierarchy as fitted. Nothing
  here fits a model, so it needs no library: \code{generate_sbm} is a
  convenience, and the statistics are language-neutral.}
\label{alg:predictive}
\begin{algorithmic}[1]
\Require $N$; $M$ draws; observed statistics $\bm{T}_{\text{obs}}$; $E$ held at
  its observed value
\For{$m = 1$ to $M$}
  \State $B \sim P(B)$;\; $\bm{n} \sim P(\bm{n} \mid B)$;\;
    $\bm{b}_{1,m} \sim P(\bm{b} \mid \bm{n})$
    \Comment{equation~\eqref{eq:pb}}
  \State $\bm{e}_{1,m} \sim P(\bm{e} \mid B, E)$, then recurse up the levels
    of \eqref{eq:nested} until one group remains
  \Comment{the whole hierarchy, not a base partition with a flat prior bolted on}
  \State $\bm{k}_m \sim P(\bm{k} \mid \bm{e}_{1,m}, \bm{b}_{1,m})$
  \State $\bm{A}_m \gets \textsc{GenerateSbm}(\bm{b}_{1,m}, \bm{e}_{1,m},
    \bm{k}_m)$, directed
  \State $\bm{T}_m \gets$ the statistics of Table~\ref{tab:stats}
\EndFor
\For{each statistic $t$}
  \State report the quantile of $t_{\text{obs}}$ within $\{t_m\}$
\EndFor
\Ensure containment, not agreement: each $t_{\text{obs}}$ should sit inside the
  bulk of its prior predictive distribution
\State \textbf{Separate comparison:} repeating steps 2--7 with a flat
  $P(\bm{e})$ in place of the recursion in step~3 isolates what the nested
  prior contributes. Run it as its own arm, labeled as such, rather than mixing the
  two.
\end{algorithmic}
\end{algorithm}

\paragraph{The interpretation.} The check is one of containment, not of
agreement. The prior does not have to generate our network; it has to generate a
universe of networks that could plausibly include it. If the observed $k$-core
profile sits far in the tail of the prior predictive distribution, the posterior
has been straining against the prior from the start -- and the 690 hub
singletons become a mechanism rather than a mystery.

% =====================================================================
\section{Stage three: posterior checks}
\label{sec:posterior}
% =====================================================================

Two kinds of checking belong here, and they answer different questions. The
first asks whether the machinery works at all. The second asks whether the
fitted model can reproduce the world.

\subsection{Three questions, not one}

Three distinct things went wrong or went unverified, and one procedure will not
settle all of them.

\paragraph{Is the inference calibrated?} Simulate parameters from the prior,
generate a network, fit, and record where the true value falls among the
posterior draws. Repeat. If prior, likelihood, and sampler are all correct,
those ranks are uniform. The summary statistic here is the number of groups,
which is discrete, so ties must be broken at random; otherwise the ranks come
out lumpy for reasons that have nothing to do with the model.

The planted-partition test already run is not this. Fitting a planted
four-block partition of 300 vertices recovered it at an adjusted Rand index of
0.96, which established that the code calls the right functions. It exercised
one favorable point in parameter space and says nothing about calibration.

Two practical constraints shape how this can be run. Under the actual prior,
$B$ is uniform on $\{1, \dots, N\}$, so most simulated networks will have
around 1{,}800 groups of two nodes each -- far from the regime we care about,
and a poor use of compute. And a single fit at full size takes 15--20 hours,
which puts $M = 100$ replications in the range of months. Calibration should
therefore be run at reduced $N$ -- a few hundred vertices -- with a short
chain at each replication. The chain is not optional: the ranks in
Algorithm~\ref{alg:sbc} need $L$ posterior draws, and the heuristic returns a
single partition. A check that uses the heuristic alone is a recovery or bias
check, not calibration, and should be named as such if it is what gets run.
What calibration at reduced $N$ validates is the inference at that scale,
which is weaker than validating it at ours and is the available option.

\begin{algorithm}[!t]
\caption{Simulation-based calibration. Answers the first question only;
  reproducibility and consensus have their own checks. Language follows
  whichever implementation is being validated.}
\label{alg:sbc}
\begin{algorithmic}[1]
\Require $M$ replications; $L$ posterior draws each; a scalar summary $f$,
  here the number of groups
\For{$m = 1$ to $M$}
  \State $\bm{b}^{*} \sim P(\bm{b})$ \Comment{the prior, not a convenient
    partition}
  \State $\bm{e}^{*}, \bm{k}^{*} \gets$ draws from their priors, as in
    Algorithm~\ref{alg:predictive}
  \State $\bm{A} \gets \textsc{GenerateSbm}(\bm{b}^{*}, \bm{e}^{*},
    \bm{k}^{*})$
  \State $\{\bm{b}^{(1)}, \dots, \bm{b}^{(L)}\} \gets$ posterior draws from
    fitting $\bm{A}$
  \State $a \gets \#\{\, l : f(\bm{b}^{(l)}) < f(\bm{b}^{*}) \,\}$;\quad
    $t \gets \#\{\, l : f(\bm{b}^{(l)}) = f(\bm{b}^{*}) \,\}$
  \State $r_m \gets a + \textsc{Binomial}(t, 1/2)$
  \Comment{$f$ is discrete; ties broken at random, or the ranks are lumpy for
    the wrong reason}
\EndFor
\Ensure $\{r_m\}$ uniform on $\{0, \dots, L\}$ if prior, likelihood, and
  sampler are correct
\State \textbf{Reading:} ranks bunched in the middle mean the posterior is too
  wide; ranks at the extremes mean it is too narrow; a consistent skew means a
  bias
\State \textbf{Scale:} at $N = 3{,}557$ a full fit is 15--20 hours, so
  $M = 100$ is months. Run at reduced $N$ -- a few hundred vertices -- with a
  short chain of $L$ draws at each replication, and say so when reporting.
  The heuristic alone yields one partition and no ranks
\end{algorithmic}
\end{algorithm}

\paragraph{Is the fit reproducible?} Two runs on identical input with the same
seed returned hierarchies 2{,}451 nats apart. Calibration will not explain
that, because it averages over replications rather than isolating a
divergence. Nor is it a mystery. \code{graph-tool}'s sweeps are parallelized with OpenMP,
and parallel reductions do not sum in a fixed order, so non-determinism under
more than one thread is documented behavior rather than a fault. The check is
a controlled repeat, and its first step is therefore
\code{openmp_set_num_threads(1)}: pin the threads, fix the seed and
\code{PYTHONHASHSEED}, run the heuristic twice in one process, and record the
partition after each internal step until the two runs differ. If they no
longer differ, the cause was the thread count and the fix is to pin it for
every reported fit. If they still differ, the remaining candidates are an
unseeded stream and iteration order over an unordered container, and the first
point of divergence names which.

\paragraph{Does the consensus recover what is there?} The step that reconciles
250 draws into one partition is the least examined part of the pipeline and the
most expensive. Algorithm~\ref{alg:sbc} never touches it, since it compares
draws against the truth directly. Check it on its own terms: simulate a network
from a known partition, fit it, run the consensus, and measure how well the
consensus recovers \emph{co-membership} -- for each pair of nodes, whether the
consensus places them together when the truth does. Co-membership rather than
labels, because labels are arbitrary and the question is whether the structure
survives the reconciliation.

\subsection{Posterior predictive checks}

Generate networks from the fitted model and compare them to the observed network
on the statistics of Table~\ref{tab:stats}. The question is now whether the
model, having seen the data, can reproduce features of it that it was not fitted
to. The $k$-core spike at coreness 74 is again the sharpest test.

\subsection{External validation}

Journals give an independent partition of the same communities, derived from
publication venue rather than from citation, and comparing the two answers the
second half of RQ7 of the specification outline, which asks whether the field's
structure follows disciplinary affiliation or cuts across it. Agreement is reported with the adjusted Rand index, computed
on modal assignments for the headline figure and across posterior draws for the
accompanying uncertainty. Two cautions belong in the reporting: the index is
sensitive to the number of groups, so the journal partition should be coarsened
to comparable granularity first; and a low value is a substantive finding about
how the field is organized, not a failure of the method.

This stage is substantive rather than diagnostic, and it is the only one of the
three that produces a result for the paper.

% =====================================================================
\section{Order of operations}
\label{sec:order}
% =====================================================================

\begin{enumerate}[leftmargin=2em, itemsep=0.4em]
\item \textbf{Settle the likelihood.} Fit twice, once with counts as
  multiplicities and once with counts as an edge covariate, and compare the
  resulting group counts and description-length curves. Everything downstream
  depends on which is in place, so nothing else can be checked first.
\item \textbf{Pushforward on depth and group sizes.} The marginal over the
  number of groups is uniform by equation~\eqref{eq:pb}; the expected
  singleton share is equation~\eqref{eq:singleton}; depth and the chance of a
  level in the elicited band are a short recursion
  (Algorithm~\ref{alg:pushforward}). None of it is simulated and none of it
  depends on $E$. Then rerun the heuristic under \code{B_max} $= 30$ and
  record the nested description length against the unconstrained fit.
\item \textbf{Decide the instrument.} If the prior cannot be brought into line
  with domain knowledge using the arguments \code{graph-tool} exposes --
  \code{B_min}, \code{B_max}, \code{clabel}, \code{pclabel} -- the blockmodel
  moves to a hand-written sampler with a stated prior. Make this decision
  explicitly and write down the reason.
\item \textbf{Prior predictive.} Generate networks from the prior; compare the
  statistics of Table~\ref{tab:stats} for containment.
\item \textbf{Calibration, reproducibility, consensus.} Three separate checks,
  set out in \S\ref{sec:posterior}. Simulation-based calibration tests the
  inference; a controlled repeat locates the non-determinism; a co-membership
  recovery test assesses the consensus step.
\item \textbf{Posterior predictive, then external validation.} The first checks
  the model against the data; the second checks the literatures against the
  journals.
\end{enumerate}

\begin{algorithm}[!t]
\caption{Settling the likelihood. Step~1 of the order above, and a
  prerequisite for Algorithm~\ref{alg:predictive}, which is the first step
  in which $E$ enters.}
\label{alg:weights}
\begin{algorithmic}[1]
\Require the community network $G$ with citation counts $\bm{w}$
\For{$\tau \in \{\texttt{eweight},\ \texttt{recs:discrete-poisson}\}$}
  \State $s_\tau \gets \textsc{MinimizeNestedBlockmodelDl}(G, \tau)$
  \State record groups per level, $\Sigma$, and $E$ under $\tau$
  \State $C_\tau \gets$ the flat description-length curve of
    Table~\ref{tab:entropy}, recomputed under $\tau$
\EndFor
\Ensure the two curves, side by side
\State \textbf{Reading:} $E$ differs by a factor of sixteen between the two
  and enters $P(\bm{e})$ directly, so a large shift in the group count is
  expected rather than surprising
\State \textbf{Note:} this is diagnostic, not decisive. No curve can say
  whether 4{,}646 citations between two communities are 4{,}646 exchangeable
  events or one intense relationship; that is a modeling judgment, argued in
  prose and recorded. The comparison shows what the judgment costs
\State \textbf{Note:} uses the 45-second heuristic, which sat 18{,}329 nats
  above the mode on the real fit. Adequate for a gross comparison of group
  counts, not for anything finer
\end{algorithmic}
\end{algorithm}

Steps one and two are inexpensive: a heuristic fit takes forty-five seconds,
and the Stage one quantities are arithmetic. Steps four through six are the substantial work, and there is no sense in
starting them before step three has settled which model they apply to.

\subsection{Where each step is implemented}

The work divides cleanly, and the dividing line is whether a step interrogates
\code{graph-tool}'s assumptions or replaces them.

The priors are written down in equations~\eqref{eq:pb} to \eqref{eq:pk}, so
nothing that only reads or samples them needs the library.
Algorithm~\ref{alg:pushforward} is arithmetic, and
Algorithm~\ref{alg:predictive} is a dozen lines of ancestral sampling in either
language. What does require Python is anything that \emph{fits}:
Algorithm~\ref{alg:weights}, the sampler inside Algorithm~\ref{alg:sbc}, and
the \code{B_max} rerun of step~2 all call \code{graph-tool}'s heuristic or its
chain, and those are built against the priors internally. The constraint is
on the fitting, not on the prior.

Anything we build ourselves should be Julia, and the usual argument for Python
evaporates at exactly the point where we stop using the package. Three reasons,
in order of force.

\emph{No probabilistic programming language handles this inference well.} A
partition is a discrete object with $B!$ label symmetry. Hamiltonian Monte
Carlo cannot move it, and Stan cannot represent discrete parameters at all.
Turing's composable Gibbs and NumPyro's \code{DiscreteHMCGibbs} can represent
them, so the door is not closed -- but single-site updates on 3{,}557 discrete
labels mix badly, the label symmetry still bites, and neither offers the
merge-and-split proposals that make this tractable. The realistic sampler is a
bespoke one: collapsed Gibbs over node reassignments, or merge-and-split moves
of the kind \code{graph-tool} uses. So the comparison is between two languages
for hand-written Markov chain Monte Carlo, not between two PPLs.

\emph{That sampler is a tight loop.} Each sweep touches 3{,}557 nodes and
36{,}105 edges, with a bookkeeping update at every reassignment. Julia runs
such loops at compiled speed as written. Python requires \code{numba},
\code{cython}, or a C extension to approach it, which means the inner loop
leaves Python anyway.

\emph{The project already has a Julia lane.} Julia 1.11.9 is in the container,
the era sweeps run there, and the Pajek readers and era-resolution code sit
beside them.

So Python's advantage here was always that someone had already written the
model. Step~3 of the order above is precisely the decision to stop relying on
that, and once it is taken the advantage is gone.

% =====================================================================
\section{What is at stake downstream}
\label{sec:downstream}
% =====================================================================

Model~I cannot be fitted until $\ell$ exists, so the blockmodel sits on the
critical path. The shape of what it returns also determines what the model can
say, and two questions about that shape are still open.

\paragraph{How $\ell$ attaches to a pair.} Every community carries a
literature label, but the slope $\alpha_{S,\ell}$ multiplies $S_{ijp}$, which
belongs to a \emph{pair}. Communities $i$ and $j$ will often sit in different
literatures, and the specification does not say what happens then. Three
readings are available and they are not equivalent: index the slope by the
citing community's literature, so $\ell = \ell_i$; index it by whether the pair
is within or across literatures, giving a two-level contrast; or carry an
effect for each of the two literatures and let them add. The first is the narrowest reading of the Model~I table in \S13.2 of the
specification outline, which lists \code{literature_i} and
\code{literature_j} as separate columns without saying how the slope consumes
them. This needs settling before any cut is chosen, because the three readings
make different demands on how many literatures there can usefully be.

\paragraph{Whether a fine cut can be estimated at all.} A grouping into 1{,}268
literatures is not collinear with the community effects -- $\theta^{K}_{i}$ and
$\psi^{K}_{j}$ are random intercepts, while $\alpha_{S,\ell}$ is a random slope
on topical proximity, and those are different objects. The real obstacle is
replication. Estimating a slope within a literature needs pairs assigned to it,
and needs those pairs to vary in $S_{ijp}$; with 2.8 communities per group at
the base, most groups would contribute a handful of pairs spanning a narrow
range of proximity, and $\sigma_{\ell}$ would be estimated almost entirely from
the prior. The quantity to report before choosing a cut is therefore simple and
has not yet been computed: for each candidate cut, the distribution across
literatures of the number of pairs and the spread of $S_{ijp}$ within them.

\paragraph{Where the candidates come from.} Four numbers have been proposed
and they have different standing. \textbf{405} is the lowest-scoring cut in
Table~\ref{tab:entropy} and is a statement about description length alone;
nobody has argued it is interpretable, and at 8.8 communities per group it
faces the replication problem above in milder form. \textbf{168} is the
coarsest cut still within 2 percent of that score. \textbf{25} is the cut at
which, by the audit of \S\ref{sec:audit}, assignment confidence stops
improving, 24 of 25 groups span more than one era, and no group is a
singleton. \textbf{8} is a reading of the hierarchy's shape; the audit prices
it at 123 percent above the best cut and finds two of its eight groups holding
2 percent of the communities between them.

The distribution of pairs across literatures, and of $S_{ijp}$ within them,
is the criterion that would separate these on grounds the model cares about,
and it is more informative than any of the arguments rehearsed above.

% =====================================================================
\section*{Sources}
% =====================================================================
\addcontentsline{toc}{section}{Sources}

\begin{itemize}[leftmargin=1.5em, itemsep=0.25em]
\item Betancourt, M. \emph{Towards a Principled Bayesian Workflow}; and
  \emph{Prior Modeling}. On the distinction between pushforward and predictive
  checks, and on eliciting priors through thresholds on interpretable
  quantities.
\item Gelman, A., Vehtari, A., Simpson, D., et al. \emph{Bayesian Workflow}
  (2020). On the sequence of checks and on weakly informative priors.
\item Talts, S., Betancourt, M., Simpson, D., Vehtari, A., \& Gelman, A.
  \emph{Validating Bayesian Inference Algorithms with Simulation-Based
  Calibration} (2018).
\item Peixoto, T. P. \emph{Hierarchical Block Structures and High-Resolution
  Model Selection in Large Networks}, Phys. Rev. X (2014). The nested
  degree-corrected blockmodel and its priors.
\item Karrer, B., \& Newman, M. E. J. \emph{Stochastic Blockmodels and Community
  Structure in Networks}, Phys. Rev. E (2011). Degree correction.
\item de Solla Price, D. J. \emph{Networks of Scientific Papers}, Science
  (1965); and \emph{A General Theory of Bibliometric and Other Cumulative
  Advantage Processes} (1976). Heavy-tailed citation distributions.
\item Small, H. \emph{Co-citation in the Scientific Literature} (1973). The
  specialty-structure tradition from which the twenty-to-thirty expectation
  comes.
\end{itemize}

\end{document}
